\documentclass[11pt,a4paper]{article}
\usepackage[utf8]{inputenc}
\usepackage[italian]{babel}
\usepackage[T1]{fontenc}
\usepackage{geometry}
\usepackage{xcolor}
\usepackage{tabularx}
\usepackage{booktabs}
\usepackage{listings}
\usepackage{tikz}
\usepackage{pgf-umlsd}
\usepackage{float}
\usepackage{hyperref}
\usepackage{amsmath}
\usepackage{enumitem}
\usepackage{tcolorbox}

\geometry{margin=2.5cm}

% ============================================
% COLORI E STILI
% ============================================

\definecolor{codebg}{RGB}{40, 42, 54}
\definecolor{codefg}{RGB}{248, 248, 242}
\definecolor{codecomment}{RGB}{98, 114, 164}
\definecolor{codekeyword}{RGB}{139, 233, 253}
\definecolor{codestring}{RGB}{255, 184, 108}
\definecolor{codefunction}{RGB}{80, 250, 123}
\definecolor{accent}{RGB}{139, 92, 246}

\lstdefinestyle{typescript}{
    language=TypeScript,
    backgroundcolor=\color{codebg},
    basicstyle=\ttfamily\small\color{codefg},
    commentstyle=\color{codecomment}\itshape,
    keywordstyle=\color{codekeyword}\bfseries,
    stringstyle=\color{codestring},
    functionstyle=\color{codefunction},
    numberstyle=\tiny\color{codecomment},
    breakatwhitespace=false,
    breaklines=true,
    captionpos=b,
    keepspaces=true,
    numbers=left,
    numbersep=5pt,
    showspaces=false,
    showstringspaces=false,
    showtabs=false,
    tabsize=2,
    frame=single,
    rulecolor=\color{accent},
    frameround=tttt,
}

% ============================================
% METADATA
% ============================================

\title{
    \textbf{Documentazione Tecnica: FlashcardList Component} \\
    \large Sistema STUDY - Modulo Flashcard
}
\author{Sistema ExtraTracker}
\date{\today}

\hypersetup{
    colorlinks=true,
    linkcolor=accent,
    filecolor=accent,
    urlcolor=accent,
    citecolor=accent,
    pdftitle={FlashcardList Component Documentation},
    pdfauthor={ExtraTracker System},
}

% ============================================
% DOCUMENTO
% ============================================

\begin{document}

\maketitle

\begin{abstract}
Questo documento fornisce un'analisi tecnica completa del componente \texttt{FlashcardList}, parte del sistema STUDY per la gestione di collezioni di flashcard. Il componente implementa un'interfaccia iOS-style con drag \& drop, supporto per viste multiple (list/grid), inserimento inline di card e integrazione con sistema di navigazione PDF. La documentazione include analisi architetturale, specifiche delle interfacce, diagrammi di flusso e pattern di interazione.
\end{abstract}

\tableofcontents
\newpage

% ============================================
% SEZIONE 1: OVERVIEW
% ============================================

\section{Overview del Componente}

\subsection{Responsabilità Principali}

Il componente \texttt{FlashcardList} rappresenta il container principale per la visualizzazione e gestione di collezioni di flashcard. Le responsabilità principali includono:

\begin{itemize}
    \item \textbf{Drag \& Drop}: Riordinamento card tramite \texttt{@dnd-kit} con feedback visivo iOS-style
    \item \textbf{Viste Multiple}: Supporto per modalità list (verticale) e grid (griglia)
    \item \textbf{Inserimento Inline}: Aggiunta card in posizioni specifiche con form inline
    \item \textbf{Gestione Stato}: Coordinamento tra stato locale e stato del deck
    \item \textbf{Filtraggio}: Supporto per visualizzazione di subset di card
    \item \textbf{Navigazione Sorgente}: Integrazione con sistema PDF per highlight card attive
    \item \textbf{Eliminazione}: Gestione eliminazione con conferma modale
\end{itemize}

\subsection{Modello Mentale: Vista Utente vs. Sistema}

\begin{tcolorbox}[colback=blue!5!white,colframe=blue!75!black,title=Modello Mentale]
\textbf{Vista Utente:}
L'utente vede una lista o griglia di card ordinabili. Può trascinare card per riordinarle, vedere pulsanti "+" tra le card per inserire nuove card, e cliccare su card per modificarle. Durante il drag, vede un overlay con la card trascinata.

\textbf{Sistema Interno:}
Il componente gestisce un DndContext con sensori configurati per prevenire drag accidentali. Mantiene stato locale ottimistico per feedback immediato, sincronizza con API solo al termine del drag. Calcola posizioni di inserimento considerando filtri attivi. Gestisce transizioni animate tra stati e previene race conditions durante operazioni asincrone.
\end{tcolorbox}

\subsection{Architettura Modulare}

Il componente è stato refactorizzato seguendo principi Clean Code in una struttura modulare:

\begin{itemize}
    \item \texttt{FlashcardList.tsx} - Componente orchestratore principale
    \item \texttt{FlashcardList.Header.tsx} - Header con conteggio e azioni
    \item \texttt{FlashcardList.InsertButton.tsx} - Pulsante inserimento inline
    \item \texttt{FlashcardList.EmptyState.tsx} - Stato vuoto
    \item \texttt{FlashcardList.DragOverlay.tsx} - Overlay durante drag
    \item \texttt{FlashcardList.constants.ts} - Costanti centralizzate
    \item \texttt{FlashcardList.types.ts} - Definizioni TypeScript
    \item \texttt{FlashcardList.utils.ts} - Funzioni pure di utilità
    \item \texttt{SortableItem.tsx} - Wrapper per card sortable
\end{itemize}

% ============================================
% SEZIONE 2: INTERFACCIA E PROPS
% ============================================

\section{Interfaccia e Props (Input)}

\subsection{Tabella Props Completa}

\begin{table}[H]
\centering
\begin{tabularx}{\textwidth}{|l|X|c|c|}
\hline
\textbf{Prop} & \textbf{Descrizione} & \textbf{Tipo} & \textbf{Req.} \\
\hline
\hline
\texttt{deck} & Oggetto Deck contenente le card e metadata. Rappresenta l'entità di dominio principale. & \texttt{Deck} & \checkmark \\
\hline
\texttt{onAddCard} & Callback per aggiunta nuova card. Invocato quando utente clicca pulsante "Aggiungi" nell'header. & \texttt{() => void} & \checkmark \\
\hline
\texttt{onUpdate} & Callback asincrono per aggiornamento contenuto card. Riceve \texttt{(cardId, front, back)} e restituisce \texttt{Promise<void>}. & \texttt{(cardId: string, front: string, back: string) => Promise<void>} & \checkmark \\
\hline
\texttt{onCardClick} & Callback opzionale invocato quando utente clicca su una card. Riceve \texttt{Card} completa. & \texttt{(card: Card) => void} & \textcolor{gray}{Opz.} \\
\hline
\texttt{onDelete} & Callback opzionale per eliminazione card. Se assente, usa \texttt{studyService.deleteCard} direttamente. & \texttt{(cardId: string) => void} & \textcolor{gray}{Opz.} \\
\hline
\texttt{showHeader} & Flag booleano per mostrare header con conteggio e azioni. Default: \texttt{false}. & \texttt{boolean} & \textcolor{gray}{Opz.} \\
\hline
\texttt{onNavigateBack} & Callback opzionale per navigazione indietro. Se presente, mostra pulsante "Indietro" nell'header. & \texttt{() => void} & \textcolor{gray}{Opz.} \\
\hline
\texttt{onDeckUpdate} & Callback opzionale invocato quando deck viene aggiornato (dopo reorder, delete, insert). Riceve \texttt{Deck} aggiornato per refresh parent. & \texttt{(updatedDeck: Deck) => void} & \textcolor{gray}{Opz.} \\
\hline
\texttt{filteredCards} & Array opzionale di card filtrate. Se fornito, usa questo invece di \texttt{deck.cards}. Utile per ricerca/filtri. & \texttt{Card[]} & \textcolor{gray}{Opz.} \\
\hline
\texttt{viewMode} & Modalità visualizzazione: \texttt{'list'} (verticale) o \texttt{'grid'} (griglia). Default: \texttt{'list'}. & \texttt{'list' | 'grid'} & \textcolor{gray}{Opz.} \\
\hline
\texttt{onShowSource} & Callback opzionale per navigazione alla fonte PDF. Passato a \texttt{FlashcardItem}. & \texttt{(card: Card) => void} & \textcolor{gray}{Opz.} \\
\hline
\texttt{activeSourceCardId} & ID opzionale della card attualmente evidenziata come "source active" (navigazione PDF). & \texttt{string | null} & \textcolor{gray}{Opz.} \\
\hline
\end{tabularx}
\caption{Tabella completa delle props di FlashcardList}
\label{tab:props}
\end{table}

\subsection{Flusso di Dati dal Parent}

Il componente riceve dati attraverso un pattern unidirezionale con supporto per filtri:

\begin{enumerate}
    \item \textbf{Props Immutabili}: Il parent passa \texttt{deck} come prop immutabile
    \item \textbf{Filtri Opzionali}: Se \texttt{filteredCards} è fornito, usa quello invece di \texttt{deck.cards}
    \item \textbf{State Locale}: Il componente mantiene \texttt{items} locale per gestire reordering ottimistico
    \item \textbf{Sincronizzazione}: \texttt{useEffect} sincronizza \texttt{items} quando \texttt{deck.cards} o \texttt{filteredCards} cambiano (solo se non in drag)
    \item \textbf{Callback Upstream}: Modifiche vengono propagate al parent tramite \texttt{onDeckUpdate}
\end{enumerate}

% ============================================
% SEZIONE 3: GESTIONE STATO E LOGICA
% ============================================

\section{Gestione dello Stato e Logica Interna}

\subsection{State Management}

Il componente utilizza sette hook di stato:

\begin{lstlisting}[style=typescript, caption=State Hooks]
const [items, setItems] = useState<Card[]>(() => {
    return filteredCards || deck.cards || [];
});
const [activeId, setActiveId] = useState<string | null>(null);
const [insertingIndex, setInsertingIndex] = useState<number | null>(null);
const [isDeleteModalOpen, setIsDeleteModalOpen] = useState(false);
const [cardToDelete, setCardToDelete] = useState<Card | null>(null);
const [isDeleting, setIsDeleting] = useState(false);
const [hoveredInsertPosition, setHoveredInsertPosition] = useState<number | null>(null);
\end{lstlisting}

\subsubsection{Analisi Dettagliata degli State}

\begin{description}
    \item[\texttt{items}] \textbf{Tipo}: \texttt{Card[]}. \textbf{Scopo}: Stato locale delle card visualizzate. Permette reordering ottimistico senza attendere API. \textbf{Inizializzazione}: Usa \texttt{filteredCards} se disponibile, altrimenti \texttt{deck.cards}. \textbf{Sincronizzazione}: Aggiornato da props solo quando non in drag (\texttt{!activeId}).
    
    \item[\texttt{activeId}] \textbf{Tipo}: \texttt{string | null}. \textbf{Scopo}: ID della card attualmente in drag. \textbf{Lifecycle}: Impostato a \texttt{card.id} in \texttt{handleDragStart}, resettato a \texttt{null} in \texttt{handleDragEnd}. \textbf{Utilizzo}: Previene sincronizzazione durante drag e mostra DragOverlay.
    
    \item[\texttt{insertingIndex}] \textbf{Tipo}: \texttt{number | null}. \textbf{Scopo}: Indice dove mostrare form inline per inserimento. \textbf{Valori}: \texttt{null} (nessun form), \texttt{0..items.length} (posizione inserimento). \textbf{Transizioni}: Impostato in \texttt{handleInsertCard}, resettato in \texttt{handleCancelInlineForm} o dopo salvataggio.
    
    \item[\texttt{isDeleteModalOpen}] \textbf{Tipo}: \texttt{boolean}. \textbf{Scopo}: Flag per visibilità modale conferma eliminazione. \textbf{Transizioni}: \texttt{true} in \texttt{handleDeleteClick}, \texttt{false} in \texttt{handleCloseDeleteModal} o dopo eliminazione.
    
    \item[\texttt{cardToDelete}] \textbf{Tipo}: \texttt{Card | null}. \textbf{Scopo}: Card selezionata per eliminazione. \textbf{Utilizzo}: Passata a \texttt{DeleteCardModal} per visualizzazione.
    
    \item[\texttt{isDeleting}] \textbf{Tipo}: \texttt{boolean}. \textbf{Scopo}: Flag per prevenire chiusura modale durante eliminazione. \textbf{Lifecycle}: \texttt{true} durante \texttt{handleConfirmDelete}, \texttt{false} nel \texttt{finally} block.
    
    \item[\texttt{hoveredInsertPosition}] \textbf{Tipo}: \texttt{number | null}. \textbf{Scopo}: Posizione hover per mostrare pulsante inserimento. \textbf{Utilizzo}: Gestito da \texttt{onMouseEnter/Leave} sui container inserimento.
\end{description}

\subsection{Effects e Sincronizzazione}

\subsubsection{useEffect per Sincronizzazione Items}

\begin{lstlisting}[style=typescript, caption=Sincronizzazione State]
useEffect(() => {
    const newCards = filteredCards || deck.cards || [];
    
    if (haveCardIdsChanged(items, newCards) && !activeId) {
        setItems(newCards);
    }
}, [deck.cards, filteredCards, activeId, items]);
\end{lstlisting}

\textbf{Logica}: L'effect sincronizza \texttt{items} locale con props \textbf{solo quando}:
\begin{enumerate}
    \item Gli ID delle card sono cambiati (non solo riordinati) - verificato da \texttt{haveCardIdsChanged()}
    \item Non si è in modalità drag (\texttt{!activeId})
\end{enumerate}

\textbf{Protezione Race Condition}: Il check \texttt{!activeId} previene sovrascrittura durante drag, preservando reordering ottimistico.

\subsection{Sensori Drag \& Drop}

\subsubsection{Configurazione Sensori}

\begin{lstlisting}[style=typescript, caption=Configurazione Sensori]
const sensors = useSensors(
    useSensor(PointerSensor, {
        activationConstraint: {
            distance: DRAG_AND_DROP.activationDistance,
        },
    }),
    useSensor(TouchSensor, {
        activationConstraint: {
            delay: DRAG_AND_DROP.touchDelay,
            tolerance: DRAG_AND_DROP.activationDistance,
        },
    }),
    useSensor(KeyboardSensor)
);
\end{lstlisting}

\textbf{Configurazione}:
\begin{itemize}
    \item \textbf{PointerSensor}: Richiede movimento di 5px prima di attivare drag (previene click accidentali)
    \item \textbf{TouchSensor}: Delay di 200ms per touch (previene scroll accidentale)
    \item \textbf{KeyboardSensor}: Supporto accessibilità per drag da tastiera
\end{itemize}

\subsection{Computed Values}

\subsubsection{Card Attiva}

\begin{lstlisting}[style=typescript, caption=Card Attiva Computed]
const activeCard = useMemo(() => {
    return activeId ? findCardById(items, activeId) : null;
}, [activeId, items]);
\end{lstlisting}

\textbf{Utilizzo}: Card visualizzata nel DragOverlay durante drag. \textbf{Memoization}: Previene ricerca ad ogni render.

\subsection{Event Handlers e Callbacks}

\subsubsection{handleDragStart}

\begin{lstlisting}[style=typescript, caption=Handler Inizio Drag]
const handleDragStart = useCallback((event: DragStartEvent) => {
    setActiveId(event.active.id as string);
}, []);
\end{lstlisting}

\textbf{Comportamento}: 
\begin{enumerate}
    \item Traccia quale card è in drag
    \item Previene sincronizzazione state durante drag
    \item Attiva DragOverlay
\end{enumerate}

\subsubsection{handleDragEnd}

\begin{lstlisting}[style=typescript, caption=Handler Fine Drag]
const handleDragEnd = useCallback(async (event: DragEndEvent) => {
    const { active, over } = event;
    setActiveId(null);
    
    if (!over || active.id === over.id) return;
    
    const oldIndex = items.findIndex((card) => card.id === active.id);
    const newIndex = items.findIndex((card) => card.id === over.id);
    
    if (oldIndex === -1 || newIndex === -1) return;
    
    try {
        const { finalOrder, newOrder } = calculateReorderResult(...);
        setItems(newOrder); // Optimistic update
        
        const updatedDeck = await studyService.reorderCards(deck.id, finalOrder);
        if (onDeckUpdate) onDeckUpdate(updatedDeck);
        
        emitToast.success(TEXT_CONTENT.toast.reorderSuccess);
    } catch (error) {
        // Revert on error
        setItems(filteredCards || deck.cards || []);
        emitToast.error(errorMessage);
    }
}, [items, deck.id, deck.cards, filteredCards, onDeckUpdate]);
\end{lstlisting}

\textbf{Flusso di Esecuzione}:
\begin{enumerate}
    \item \textbf{Reset State}: \texttt{setActiveId(null)} per permettere sincronizzazione
    \item \textbf{Guard Clauses}: Verifica \texttt{over} presente e posizione cambiata
    \item \textbf{Calcolo Indici}: Trova posizioni vecchia e nuova
    \item \textbf{Optimistic Update}: Aggiorna \texttt{items} immediatamente per feedback
    \item \textbf{API Call}: Chiama \texttt{studyService.reorderCards} con ordine finale
    \item \textbf{Success Path}: Aggiorna parent e mostra toast
    \item \textbf{Error Path}: Reverta state locale e mostra errore
\end{enumerate}

\textbf{Calcolo Ordine Finale}: La funzione \texttt{calculateReorderResult()} gestisce il caso di card filtrate, mappando l'ordine visibile all'ordine completo del deck.

\subsubsection{handleSaveInlineForm}

\begin{lstlisting}[style=typescript, caption=Handler Salvataggio Form Inline]
const handleSaveInlineForm = useCallback(async (front: string, back: string) => {
    if (insertingIndex === null) return;
    
    try {
        const insertPosition = calculateInsertPosition(
            insertingIndex,
            filteredCards,
            deck.cards || []
        );
        
        const updatedDeck = await studyService.addCardAtPosition(deck.id, {
            front: front.trim(),
            back: back.trim(),
            position: insertPosition,
        });
        
        if (onDeckUpdate) onDeckUpdate(updatedDeck);
        
        emitToast.success(TEXT_CONTENT.toast.addSuccess(insertPosition));
        setInsertingIndex(null);
    } catch (error) {
        emitToast.error(errorMessage);
    }
}, [insertingIndex, deck.id, deck.cards, filteredCards, onDeckUpdate]);
\end{lstlisting}

\textbf{Logica Posizione Inserimento}: La funzione \texttt{calculateInsertPosition()} mappa la posizione nel view filtrato alla posizione reale nel deck completo:
\begin{itemize}
    \item Se \texttt{position === 0}: Inserisce all'inizio del deck
    \item Se \texttt{position >= filteredCards.length}: Inserisce alla fine del deck
    \item Altrimenti: Trova posizione del card precedente nel deck completo e inserisce dopo
\end{itemize}

% ============================================
% SEZIONE 4: INTERAZIONI E COMUNICAZIONE
% ============================================

\section{Interazioni e Comunicazione (I/O)}

\subsection{Comunicazione Upstream (Parent)}

Il componente comunica con il parent attraverso callbacks:

\begin{description}
    \item[\texttt{onAddCard()}] \textbf{Trigger}: Click pulsante "Aggiungi" nell'header. \textbf{Payload}: Nessuno. \textbf{Contratto}: Sincrono, gestione aggiunta delegata al parent.
    
    \item[\texttt{onUpdate(cardId, front, back)}] \textbf{Trigger}: Salvataggio modifiche da \texttt{FlashcardItem}. \textbf{Payload}: ID card e contenuti validati. \textbf{Contratto}: Promise che risolve quando persistenza completata.
    
    \item[\texttt{onDeckUpdate(updatedDeck)}] \textbf{Trigger}: Dopo reorder, delete, o insert completati. \textbf{Payload}: Deck aggiornato completo. \textbf{Contratto}: Sincrono, permette refresh parent.
    
    \item[\texttt{onCardClick(card)}] \textbf{Trigger}: Click su card (view mode). \textbf{Payload}: Oggetto card completo. \textbf{Contratto}: Sincrono, nessun ritorno.
    
    \item[\texttt{onDelete(cardId)}] \textbf{Trigger}: Conferma eliminazione da modale. \textbf{Payload}: Solo ID. \textbf{Contratto}: Promise, gestione eliminazione delegata al parent se fornito.
    
    \item[\texttt{onShowSource(card)}] \textbf{Trigger}: Click pulsante "Show Source" su card. \textbf{Payload}: Card completa. \textbf{Contratto}: Sincrono, navigazione PDF gestita dal parent.
    
    \item[\texttt{onNavigateBack()}] \textbf{Trigger}: Click pulsante "Indietro" nell'header. \textbf{Payload}: Nessuno. \textbf{Contratto}: Sincrono, navigazione gestita dal parent.
\end{description}

\subsection{Comunicazione Downstream (Children)}

Il componente passa dati a sub-componenti:

\begin{description}
    \item[\texttt{Header}] Riceve: \texttt{cardCount}, \texttt{showBackButton}, \texttt{showAddButton}, callbacks. Passa: Nessun dato, solo configurazione.
    
    \item[\texttt{SortableItem}] Riceve: \texttt{card}, \texttt{index}, \texttt{totalCards}, \texttt{onUpdate}, \texttt{onClick}, \texttt{onDelete}, \texttt{viewMode}, \texttt{onShowSource}, \texttt{isSourceActive}. Passa: Props per \texttt{FlashcardItem}.
    
    \item[\texttt{FlashcardInlineForm}] Riceve: \texttt{onSave}, \texttt{onCancel}, \texttt{position}. Passa: Handlers per salvataggio/annullamento.
    
    \item[\texttt{InsertButton}] Riceve: \texttt{position}, \texttt{onInsert}, \texttt{isVisible}. Passa: Configurazione visibilità e handler click.
    
    \item[\texttt{DragOverlayContent}] Riceve: \texttt{card}, \texttt{viewMode}, \texttt{items}, tutti i callbacks. Passa: Props per \texttt{FlashcardItem} con styling overlay.
    
    \item[\texttt{DeleteCardModal}] Riceve: \texttt{isOpen}, \texttt{onClose}, \texttt{onConfirm}, \texttt{card}, \texttt{isDeleting}. Passa: Stato modale e card da eliminare.
\end{description}

\subsection{Interazioni Esterne}

\begin{description}
    \item[\texttt{studyService}] \textbf{Modulo}: \texttt{services/studyService}. \textbf{Metodi Usati}:
    \begin{itemize}
        \item \texttt{reorderCards(deckId, order)}: Riordina card nel deck
        \item \texttt{addCardAtPosition(deckId, cardData)}: Aggiunge card in posizione specifica
        \item \texttt{deleteCard(deckId, cardId)}: Elimina card (se \texttt{onDelete} non fornito)
    \end{itemize}
    
    \item[\texttt{emitToast}] \textbf{Modulo}: \texttt{shared/components/toast}. \textbf{Uso}: Feedback utente per successo/errore operazioni. \textbf{Pattern}: Toast notifications non bloccanti.
    
    \item[\texttt{@dnd-kit}] \textbf{Libreria}: Drag \& drop library. \textbf{Componenti Usati}:
    \begin{itemize}
        \item \texttt{DndContext}: Context provider per drag operations
        \item \texttt{SortableContext}: Context per sortable items
        \item \texttt{DragOverlay}: Overlay durante drag
        \item \texttt{useSortable}: Hook per rendere items sortable
    \end{itemize}
\end{description}

% ============================================
% SEZIONE 5: DIAGRAMMI TECNICI
% ============================================

\section{Diagrammi Tecnici}

\subsection{Diagramma di Sequenza: Flusso Drag \& Drop}

\begin{figure}[H]
\centering
\begin{tikzpicture}[
    node distance=1.5cm,
    every node/.style={align=center},
    component/.style={rectangle, draw=accent, fill=accent!10, minimum width=2.5cm, minimum height=1cm},
    action/.style={rectangle, draw=blue!50, fill=blue!5, minimum width=2cm, minimum height=0.8cm},
]

% Nodes
\node[component] (user) {User};
\node[component, right=of user] (list) {FlashcardList};
\node[component, right=of list] (dnd) {DndContext};
\node[component, below=of list] (item) {SortableItem};
\node[component, below=of dnd] (api) {studyService};

% Drag Start
\draw[->, dashed, blue] (user) -- node[above] {Start Drag} (list);
\draw[->, thick, blue] (list) -- node[above] {\texttt{handleDragStart()}} (dnd);
\draw[->, thick, blue] (list) -- node[right] {\texttt{setActiveId(cardId)}} (item);
\draw[->, thick, green] (list) -- node[below, sloped] {Show DragOverlay} (list);

% During Drag
\draw[->, dashed, blue] (user) -- node[above] {Move} (dnd);
\draw[->, thick, blue] (dnd) -- node[right] {Update Transform} (item);

% Drag End
\draw[->, dashed, blue] (user) -- node[above] {Drop} (list);
\draw[->, thick, blue] (list) -- node[above] {\texttt{handleDragEnd()}} (dnd);
\draw[->, thick, blue] (list) -- node[right] {\texttt{calculateReorderResult()}} (list);
\draw[->, thick, green] (list) -- node[right] {\texttt{setItems(newOrder)}} (list);
\draw[->, thick, red] (list) -- node[below] {\texttt{reorderCards(deckId, order)}} (api);

% Response
\draw[->, dashed, green] (api) -- node[below] {\texttt{Promise.resolve(deck)}} (list);
\draw[->, thick, green] (list) -- node[above, sloped] {\texttt{onDeckUpdate(deck)}} (list);
\draw[->, thick, green] (list) -- node[right] {\texttt{emitToast.success()}} (list);

\end{tikzpicture}
\caption{Diagramma di sequenza: Flusso drag \& drop completo}
\label{fig:sequence-drag}
\end{figure}

\subsection{Diagramma di Flusso: State Machine Drag}

\begin{figure}[H]
\centering
\begin{tikzpicture}[
    node distance=2.5cm,
    state/.style={ellipse, draw=accent, fill=accent!10, minimum width=2.5cm, minimum height=1.5cm},
    transition/.style={->, thick, blue},
    decision/.style={diamond, draw=orange, fill=orange!10, minimum width=2cm, minimum height=1.5cm},
]

% States
\node[state] (idle) {Idle\\State};
\node[state, right=of idle] (dragging) {Dragging\\State};
\node[state, below=of dragging] (syncing) {Syncing\\State};
\node[state, left=of syncing] (error) {Error\\State};

% Transitions
\draw[transition] (idle) -- node[above] {\texttt{handleDragStart()}} (dragging);
\draw[transition] (dragging) -- node[right] {\texttt{handleDragEnd()}} (syncing);
\draw[transition] (syncing) -- node[below] {\texttt{Success}} (idle);
\draw[transition] (syncing) -- node[left] {\texttt{Error}} (error);
\draw[transition] (error) -- node[below] {\texttt{Revert State}} (idle);

% Decision points
\node[decision, above=of syncing] (validate) {Valid\\Drop?};
\draw[transition] (dragging) -- (validate);
\draw[transition] (validate) -- node[left] {No} (idle);
\draw[transition] (validate) -- node[right] {Yes} (syncing);

\end{tikzpicture}
\caption{State Machine: Transizioni di stato durante drag}
\label{fig:statemachine-drag}
\end{figure}

\subsection{Diagramma Architetturale: Struttura Modulare}

\begin{figure}[H]
\centering
\begin{tikzpicture}[
    node distance=1.8cm,
    module/.style={rectangle, draw=accent, fill=accent!10, rounded corners, minimum width=3cm, minimum height=1.2cm, text centered},
    data/.style={rectangle, draw=green!50, fill=green!5, rounded corners, minimum width=2.5cm, minimum height=0.8cm},
    arrow/.style={->, thick},
]

% Main component
\node[module] (main) {FlashcardList.tsx\\Orchestrator};

% Sub-components
\node[module, above=of main] (header) {Header.tsx};
\node[module, right=of main] (sortable) {SortableItem.tsx};
\node[module, below=of main] (insert) {InsertButton.tsx};
\node[module, left=of main] (empty) {EmptyState.tsx};
\node[module, below=of sortable] (overlay) {DragOverlay.tsx};

% Supporting modules
\node[data, above=of header] (constants) {constants.ts};
\node[data, right=of constants] (types) {types.ts};
\node[data, below=of types] (utils) {utils.ts};

% Connections
\draw[arrow] (main) -- node[left] {renders} (header);
\draw[arrow] (main) -- node[above] {renders} (sortable);
\draw[arrow] (main) -- node[left] {renders} (insert);
\draw[arrow] (main) -- node[below] {renders} (empty);
\draw[arrow] (main) -- node[right] {renders} (overlay);
\draw[arrow] (constants) -- node[above] {imports} (main);
\draw[arrow] (types) -- node[above] {imports} (main);
\draw[arrow] (utils) -- node[below] {imports} (main);

\end{tikzpicture}
\caption{Architettura modulare: Dipendenze e struttura}
\label{fig:architecture}
\end{figure}

\subsection{Diagramma di Flusso: Calcolo Posizione Inserimento}

\begin{figure}[H]
\centering
\begin{tikzpicture}[
    node distance=2cm,
    process/.style={rectangle, draw=blue!50, fill=blue!5, minimum width=3cm, minimum height=1cm, text centered},
    decision/.style={diamond, draw=orange, fill=orange!10, minimum width=2.5cm, minimum height=1.5cm, text centered},
    startend/.style={ellipse, draw=green!50, fill=green!5, minimum width=2.5cm, minimum height=1cm, text centered},
]

% Nodes
\node[startend] (start) {Inizio};
\node[decision, below=of start] (check) {Filtri\\attivi?};
\node[process, right=of check] (simple) {Return\\position};
\node[decision, below=of check] (pos) {position\\== 0?};
\node[process, left=of pos] (begin) {Return 0};
\node[decision, below=of pos] (end) {position\\>= length?};
\node[process, right=of end] (last) {Return\\deck.length};
\node[process, below=of end] (calc) {Find previous\\card index};
\node[process, below=of calc] (result) {Return\\prevIndex + 1};
\node[startend, below=of result] (finish) {Fine};

% Connections
\draw[->, thick] (start) -- (check);
\draw[->, thick] (check) -- node[above] {No} (simple);
\draw[->, thick] (check) -- node[left] {Yes} (pos);
\draw[->, thick] (pos) -- node[left] {Yes} (begin);
\draw[->, thick] (pos) -- node[right] {No} (end);
\draw[->, thick] (end) -- node[above] {Yes} (last);
\draw[->, thick] (end) -- node[right] {No} (calc);
\draw[->, thick] (calc) -- (result);
\draw[->, thick] (simple) |- (finish);
\draw[->, thick] (begin) |- (finish);
\draw[->, thick] (last) |- (finish);
\draw[->, thick] (result) -- (finish);

\end{tikzpicture}
\caption{Flusso di calcolo posizione inserimento con filtri}
\label{fig:insert-position}
\end{figure}

% ============================================
% SEZIONE 6: EDGE CASES E ERROR HANDLING
% ============================================

\section{Edge Cases e Error Handling}

\subsection{Gestione Dati Mancanti}

\begin{description}
    \item[\texttt{deck} null/undefined] \textbf{Prevenzione}: TypeScript garantisce prop required. \textbf{Fallback}: Componente non renderizza se deck mancante (parent responsabilità).
    
    \item[\texttt{deck.cards} vuoto/null] \textbf{Comportamento}: Mostra \texttt{EmptyState} con messaggio informativo. \textbf{Logica}: Check \texttt{cardCount === 0} nel render.
    
    \item[\texttt{filteredCards} vuoto] \textbf{Comportamento}: Usa \texttt{deck.cards} come fallback. \textbf{Logica}: \texttt{filteredCards || deck.cards || []} in inizializzazione state.
    
    \item[Card non trovata durante drag] \textbf{Protezione}: Guard clause \texttt{if (oldIndex === -1 || newIndex === -1) return} in \texttt{handleDragEnd}. \textbf{Comportamento}: Drag viene ignorato silenziosamente.
\end{description}

\subsection{Gestione Errori API}

\begin{lstlisting}[style=typescript, caption=Error Handling Pattern]
try {
    const updatedDeck = await studyService.reorderCards(deck.id, finalOrder);
    if (onDeckUpdate) onDeckUpdate(updatedDeck);
    emitToast.success(TEXT_CONTENT.toast.reorderSuccess);
} catch (error) {
    const errorMessage = error instanceof Error 
        ? error.message 
        : TEXT_CONTENT.toast.reorderError;
    
    console.error('[FlashcardList] Errore nel riordinamento:', error);
    emitToast.error(errorMessage);
    
    // Revert local state on error
    setItems(filteredCards || deck.cards || []);
}
\end{lstlisting}

\textbf{Strategia}:
\begin{enumerate}
    \item \textbf{Type Guard}: Verifica che error sia istanza di \texttt{Error}
    \item \textbf{Fallback Message}: Usa messaggio default se error non standard
    \item \textbf{State Revert}: Ripristina state locale a valori originali
    \item \textbf{User Feedback}: Mostra toast di errore per informare utente
    \item \textbf{Logging}: Log error per debugging (console.error)
\end{enumerate}

\subsection{Concorrenza e Race Conditions}

\begin{description}
    \item[Drag durante sincronizzazione] \textbf{Problema}: Props cambiano durante drag, causando sovrascrittura state locale. \textbf{Soluzione}: Check \texttt{!activeId} in \texttt{useEffect} previene sincronizzazione durante drag.
    
    \item[Doppi drag simultanei] \textbf{Problema}: Utente inizia nuovo drag prima che precedente finisca. \textbf{Soluzione}: \texttt{@dnd-kit} gestisce internamente, solo un drag attivo alla volta.
    
    \item[Inserimento durante drag] \textbf{Problema}: Form inline aperto durante drag. \textbf{Comportamento}: Form rimane aperto, drag non interferisce (gestito da \texttt{insertingIndex} separato).
    
    \item[Eliminazione durante drag] \textbf{Problema}: Card eliminata mentre in drag. \textbf{Protezione}: \texttt{handleDragEnd} verifica esistenza card prima di processare.
\end{description}

\subsection{Gestione Filtri e Reordering}

\begin{description}
    \item[Reordering con filtri attivi] \textbf{Problema}: Ordine visibile non corrisponde a ordine completo deck. \textbf{Soluzione}: \texttt{calculateReorderResult()} mappa ordine filtrato a ordine completo, preservando posizioni card nascoste.
    
    \item[Inserimento con filtri attivi] \textbf{Problema}: Posizione nel view filtrato non corrisponde a posizione reale. \textbf{Soluzione}: \texttt{calculateInsertPosition()} calcola posizione reale considerando card filtrate.
    
    \item[Sincronizzazione dopo filtri] \textbf{Problema}: Filtri cambiano mentre drag in corso. \textbf{Comportamento}: Drag continua con card originali, sincronizzazione avviene dopo \texttt{handleDragEnd}.
\end{description}

\subsection{Accessibilità e UX Edge Cases}

\begin{description}
    \item[Keyboard Navigation] \textbf{Supporto}: \texttt{KeyboardSensor} abilita drag da tastiera. \textbf{Limiti}: Navigazione tra card tramite Tab standard, drag attivato con Space/Enter.
    
    \item[Screen Readers] \textbf{Supporto}: \texttt{aria-label} su tutti i pulsanti interattivi. \textbf{Contenuto}: Labels descrittivi per azioni (Insert, Delete, Reorder).
    
    \item[Touch Devices] \textbf{Supporto}: \texttt{TouchSensor} con delay per prevenire conflitti con scroll. \textbf{Area Touch}: Card intere sono draggable, non solo handle.
    
    \item[Performance con molte card] \textbf{Problema}: Render di centinaia di card può essere lento. \textbf{Ottimizzazione}: \texttt{React.memo} su \texttt{SortableItem}, virtualizzazione non implementata (possibile miglioramento futuro).
\end{description}

% ============================================
% SEZIONE 7: SUB-COMPONENTI
% ============================================

\section{Sub-Componenti}

\subsection{Header Component}

\subsubsection{Responsabilità}

Il componente \texttt{Header} gestisce la visualizzazione del conteggio card e dei pulsanti di azione (Indietro, Aggiungi).

\subsubsection{Props}

\begin{tabularx}{\textwidth}{|l|X|c|}
\hline
\textbf{Prop} & \textbf{Descrizione} & \textbf{Tipo} \\
\hline
\texttt{cardCount} & Numero totale di card & \texttt{number} \\
\hline
\texttt{showBackButton} & Flag per mostrare pulsante Indietro & \texttt{boolean} \\
\hline
\texttt{showAddButton} & Flag per mostrare pulsante Aggiungi & \texttt{boolean} \\
\hline
\texttt{onNavigateBack} & Callback navigazione indietro & \texttt{() => void} \\
\hline
\texttt{onAddCard} & Callback aggiunta card & \texttt{() => void} \\
\hline
\end{tabularx}

\subsubsection{Logica di Rendering}

Il componente mostra badge con conteggio e pulsanti condizionali basati su flags:

\begin{lstlisting}[style=typescript, caption=Render Condizionale]
{showBackButton && onNavigateBack && (
    <button onClick={onNavigateBack}>...</button>
)}
{showAddButton && (
    <button onClick={onAddCard}>...</button>
)}
\end{lstlisting}

\subsection{InsertButton Component}

\subsubsection{Responsabilità}

Il componente \texttt{InsertButton} gestisce la visualizzazione del pulsante "+" per inserimento inline, con animazioni e feedback visivo.

\subsubsection{Props}

\begin{tabularx}{\textwidth}{|l|X|c|}
\hline
\textbf{Prop} & \textbf{Descrizione} & \textbf{Tipo} \\
\hline
\texttt{position} & Posizione inserimento & \texttt{number} \\
\hline
\texttt{onInsert} & Callback click pulsante & \texttt{(position: number) => void} \\
\hline
\texttt{isVisible} & Flag visibilità pulsante & \texttt{boolean} \\
\hline
\end{tabularx}

\subsubsection{Animazioni}

Il componente utilizza \texttt{framer-motion} per:
\begin{itemize}
    \item \textbf{Entrata/Uscita}: Fade e scale quando \texttt{isVisible} cambia
    \item \textbf{Linea Animata}: Pulsazione orizzontale continua
    \item \textbf{Icona Animata}: Pulsazione con glow effect
\end{itemize}

\subsection{SortableItem Component}

\subsubsection{Responsabilità}

Il componente \texttt{SortableItem} è un wrapper che rende \texttt{FlashcardItem} draggable usando \texttt{@dnd-kit}'s \texttt{useSortable} hook.

\subsubsection{Logica Drag}

\begin{lstlisting}[style=typescript, caption=useSortable Hook]
const {
    attributes,
    listeners,
    setNodeRef,
    transform,
    transition,
    isDragging,
} = useSortable({ id: card.id });

const style = {
    transform: CSS.Transform.toString(transform),
    transition: isDragging ? 'all 200ms ease-out' : transition,
};
\end{lstlisting}

\textbf{Comportamento}:
\begin{itemize}
    \item \textbf{Transform}: Applicato automaticamente da \texttt{@dnd-kit} durante drag
    \item \textbf{Ghost Effect}: Opacità ridotta e scale quando \texttt{isDragging === true}
    \item \textbf{View Modes}: Layout diverso per list (con handle) vs grid (semplice)
\end{itemize}

\subsection{DragOverlayContent Component}

\subsubsection{Responsabilità}

Il componente \texttt{DragOverlayContent} visualizza la card in drag con styling speciale (scale, rotate, shadow).

\subsubsection{Styling Overlay}

\begin{lstlisting}[style=typescript, caption=Styling Drag Overlay]
const overlayStyles = {
    scale: 'scale-110',
    rotate: 'rotate-3',
    shadow: 'shadow-2xl',
    cursor: 'cursor-grabbing',
    boxShadow: '0 25px 50px -12px rgba(139, 92, 246, 0.5)',
    filter: 'drop-shadow(0 30px 60px -15px rgba(0, 0, 0, 0.4))',
};
\end{lstlisting}

\textbf{Effetti Visivi}:
\begin{itemize}
    \item \textbf{Scale}: 110\% per enfatizzare card in drag
    \item \textbf{Rotate}: 3° per effetto dinamico
    \item \textbf{Shadow}: Ombra pronunciata con colore accent
    \item \textbf{Cursor}: Cambio a "grabbing" durante drag
\end{itemize}

% ============================================
% SEZIONE 8: PERFORMANCE E OTTIMIZZAZIONI
% ============================================

\section{Performance e Ottimizzazioni}

\subsection{Optimistic Updates}

Il componente implementa optimistic updates per feedback immediato:

\begin{lstlisting}[style=typescript, caption=Optimistic Update Pattern]
// Update local state immediately
setItems(newOrder);

// Then sync with API
const updatedDeck = await studyService.reorderCards(deck.id, finalOrder);
\end{lstlisting}

\textbf{Benefici}: Utente vede cambiamento immediato, API call non blocca UI.

\textbf{Revert Strategy}: In caso di errore, state viene ripristinato a valori originali.

\subsection{useMemo per Computed Values}

Valori computati sono memoizzati:

\begin{lstlisting}[style=typescript, caption=Memoization]
const activeCard = useMemo(() => {
    return activeId ? findCardById(items, activeId) : null;
}, [activeId, items]);
\end{lstlisting}

\textbf{Benefici}: Previene ricerca card ad ogni render, solo quando \texttt{activeId} o \texttt{items} cambiano.

\subsection{useCallback per Event Handlers}

Tutti gli event handlers utilizzano \texttt{useCallback}:

\begin{itemize}
    \item \texttt{handleDragStart}: Dipende da nessuna variabile (stabile)
    \item \texttt{handleDragEnd}: Dipende da \texttt{items}, \texttt{deck.id}, \texttt{deck.cards}, \texttt{filteredCards}, \texttt{onDeckUpdate}
    \item \texttt{handleSaveInlineForm}: Dipende da \texttt{insertingIndex}, \texttt{deck.id}, \texttt{deck.cards}, \texttt{filteredCards}, \texttt{onDeckUpdate}
\end{itemize}

\textbf{Impatto}: Previene ricreazione funzioni ad ogni render, riducendo re-render di children.

\subsection{Prevenzione Re-render Inutili}

\begin{description}
    \item[Sincronizzazione Condizionale] \textbf{Pattern}: \texttt{useEffect} sincronizza solo quando necessario (card IDs cambiati, non in drag). \textbf{Beneficio}: Evita re-render durante drag operations.
    
    \item[React.memo su Children] \textbf{Pattern}: \texttt{SortableItem} e sub-componenti potrebbero beneficiare di \texttt{React.memo}. \textbf{Stato Attuale}: Non implementato, possibile ottimizzazione futura.
\end{description}

% ============================================
% CONCLUSIONI
% ============================================

\section{Conclusioni}

Il componente \texttt{FlashcardList} rappresenta un'implementazione robusta e production-ready di un'interfaccia drag \& drop per la gestione di collezioni di flashcard. L'architettura modulare, la gestione ottimistica dello stato e l'integrazione con \texttt{@dnd-kit} garantiscono performance, manutenibilità e UX fluida.

\textbf{Punti di Forza}:
\begin{itemize}
    \item Architettura modulare con separazione concerns
    \item Drag \& drop fluido con feedback visivo iOS-style
    \item Gestione filtri complessa con mapping corretto posizioni
    \item Optimistic updates per feedback immediato
    \item Error handling robusto con revert automatico
    \item Supporto accessibilità con keyboard navigation
\end{itemize}

\textbf{Aree di Miglioramento Future}:
\begin{itemize}
    \item Implementazione virtualizzazione per liste molto lunghe
    \item Aggiunta unit tests per funzioni pure di utilità
    \item Implementazione undo/redo per operazioni drag
    \item Ottimizzazione animazioni per dispositivi low-end
    \item Aggiunta React.memo su sub-componenti per performance
    \item Supporto multi-select per operazioni batch
\end{itemize}

\end{document}
